\section{Quantized deployment validation}
\label{app:awq}

The AWQ study extends the closed-loop validation contract to one quantized
serving configuration. It is a deployment-context validation, rather than a
quantization-only causal comparison: the quantized weights, kernel, loading
path, and runtime are all part of the new configuration. The policies,
workloads, deadlines, accounting prices, and guard rule remain those fixed in
the main study.

\paragraph{Fixed configuration.}
We use Qwen2.5-7B-Instruct-AWQ with vLLM 0.9.0 and the AWQ-Marlin backend on
one host with two A100 80GB PCIe GPUs. Each GPU runs an independent TP=1
instance with FP16 activations, four local slots, a 4096-token context, eager
execution, chunked prefill, and no prefix cache. The preceding precheck
measured 19 service shapes (three repeats per shape on each GPU), 20 loaded
scale-out cycles, and separately measured setup and release. It froze three
CPU predictors that share the new service and guard parameters and differ only
in startup: AWQ-Hist reuses the historical 69.219 s proxy, AWQ-30 uses 30 s,
and AWQ-Local uses the measured 25.810 s median. Only AWQ-Local drives the
physical controller; the other two are counterfactual CPU replays.

\paragraph{Physical matrix.}
The formal matrix has two locked windows, recovery and steady, and three
policies: all2, H (the hysteresis policy), and H+guard. Each cell is repeated
twice for 12 physical runs. The four all2 cells are capacity anchors. The
remaining eight cells are compared with the frozen AWQ-Local prediction made
before execution. A run is technically valid only after worker identity,
request terminals, lifecycle cleanup, and the resource ledger pass independent
CPU checks. No technical retry was used.

\begin{table}[t]
\centering
\small
\caption{AWQ-Marlin physical results, averaged over two repeats. The prediction
and physical columns report P1+ passes over the two repeats; offline counts use
120 arrivals per repeat. Occupancy (GPU seconds) and cost (scenario USD) are physical means.}
\label{tab:awq-results}
\setlength{\tabcolsep}{2pt}
\par\vspace{4pt}
\begin{tabular*}{\linewidth}{@{\extracolsep{\fill}}llrrrrr@{}}
\toprule
Window & Policy & Pred. & Phys. & Offline & GPU sec. & Cost\\
\midrule
Recovery & H & 2/2 & 2/2 & 120/120 & 788.505 & 0.9536\\
Recovery & H+guard & 2/2 & 2/2 & 120/120 & 788.623 & 0.9537\\
Steady & H & 0/2 & 0/2 & 119/120 & 925.882 & 1.2698\\
Steady & H+guard & 2/2 & 2/2 & 120/120 & 977.267 & 1.3137\\
\bottomrule
\end{tabular*}
\end{table}

\begin{table}[t]
\centering
\small
\caption{Frozen AWQ predictions versus eight policy runs, excluding four
capacity anchors. FA/FR count false accepts/rejects. Occupancy MAE is in GPU
seconds; cost MAE is in scenario USD. Four policy/window predictions per model
are each reused for two physical repeats.}
\label{tab:awq-predictors}
\setlength{\tabcolsep}{2pt}
\par\vspace{4pt}
\begin{tabular*}{\linewidth}{@{\extracolsep{\fill}}lrrrrr@{}}
\toprule
Predictor & Correct & FA & FR & Occ. MAE & Cost MAE\\
\midrule
AWQ-Hist & 6/8 & 0 & 2 & 868.054 & 0.83361\\
AWQ-30 & 8/8 & 0 & 0 & 8.495 & 0.02247\\
AWQ-Local & 8/8 & 0 & 0 & 2.745 & 0.00436\\
\bottomrule
\end{tabular*}
\end{table}

\paragraph{Results and interpretation.}
All four all2 anchors pass P1+, and all eight policy cells are technically
valid. AWQ-Local predicts all eight policy outcomes correctly: the six feasible
cells pass and the two H/steady cells fail in both repeats. Those two physical
failures are scientific negative results, not technical failures: every online
request is on time, while 119 of 120 offline requests meet their deadlines and
one remains unfinished at the 2100 s cutoff. H+guard completes the same steady
window in both repeats. The frozen policy target matches the physical target at
all 2,100 ticks in every policy cell. Across the eight policy cells, mean
absolute error is 2.745 GPU seconds for occupied resource and 0.00436 scenario
USD for total cost; the corresponding mean trajectory L1 error is 3.143 GPU
seconds.

Table~\ref{tab:awq-predictors} compares all three frozen predictors against the
same physical runs. AWQ-Hist falsely rejects both feasible H+guard steady
repeats; AWQ-30 matches all eight outcomes with 3.09 times AWQ-Local's occupancy
MAE. Thus historical startup reuse can change feasibility judgments even with
shared local service estimates. All three correctly predict the H/steady failures.

The AWQ archive retains the calibration records, frozen predictions, raw
controller and worker events, per-cell request terminals, lifecycle ledgers,
and independent CPU audit files. The reduced predictions and ledgers are recomputed by
\path{artifact/awq_formal/verify.py}. These files are separate from the public
197-run corpus counts in Appendix~A; the AWQ study should be counted as an
additional deployment-context extension when the archive is made public.
